\subsection{The diagonal pairs with a nonprimitive frame span}

The partner count used for a generic term fixes its full integral
orthogonal lattice. A nonprimitive pair of frames can have different
integral orthogonal lattices even when the rational orthogonal spaces
agree modulo the flag. Retaining the two indices in the common
congruence lattice gives the required bound at corank one.

\begin{auxlemma}[Corank one with a nonprimitive frame span]\label{lem:nonprimitive-diagonal}
Fix a marked word \(u\) on the row and use the moment order of
Lemma~\ref{lem:row-moment-budget}. In the matching ordered pairs on
the block \((u,u)\), let \(\mathcal S\) be any subset whose two
single-term Gram invariants are positive and whose common rank is
\(k-1\). With the weights, finite norm and coefficient norms of
Lemma~\ref{lem:equal-word-generic-corank-one},
\begin{equation}\label{eq:nonprimitive-diagonal}
  \Gamma_1\sum_{(t,t')\in\mathcal S}w_{t,t'}|a_t a_{t'}|
       \le \exp(C_N m^5f\log b)\,\theta_{uu}.
\end{equation}
Here \(C_N\) depends only on the fixed counting, level and local
amplitude constants and can be chosen before \(C_\Gamma\).
In particular the bound applies to the pairs in the nongeneric
term group on the diagonal.
\end{auxlemma}

\begin{proof}
\proofstep{The indices of one marked pair of frames.}
Write \(R=\mathcal O_E\), \(g=g_0g_1\), \(p=r+1\), and
\(N=T_u^f\). In an integral unimodular row lattice \(\Lambda\)
denote the fixed flag by \((X,z)\) and its two partners by
\((X_j,z_j)\), \(j=1,2\). Here \(X,X_j\) are the isotropic spaces
and \(z,z_j\) the distinguished vectors. Put
\[
 \begin{aligned}
 S_j&=(\Lambda\cap X)+Rz+(\Lambda\cap X_j)+Rz_j,&
 W_j&=E S_j,\\
 B_j&=(\Lambda\cap W_j)/S_j,&
 U_j&=\Lambda\cap W_j^\perp,\qquad
 \beta_j=|B_j|/g^{mf}.
 \end{aligned}
\]
The two augmented flag spaces in \(W_j\) form a direct sum by
positivity of the single-term Gram invariant. Let \(G_j\) be the
normalized determinant \(G_Z\) of the corresponding reduced label.
The determinant calculation in Proposition~\ref{prop:flag-pair-counts}
before taking the primitive hull gives \(d(S_j)=g^{mf}G_j\).
The saturated lattice \(\Lambda\cap W_j\) and \(U_j\) have
opposite discriminant pairings of equal order by unimodularity.
Thus in the determinant convention \(d(L)=\sqrt{|L^\#/L|}\),
\[
   d(U_j)=d(\Lambda\cap W_j)
          =\frac{d(S_j)}{|B_j|}=\frac{G_j}{\beta_j}.
\]

At primes over \(g\), the marked flags have unit integral bases and
their level comparison is upper unitriangular. Dividing the \(p\)
corresponding frame differences by \(g\) gives integral classes
in \(B_j[g]\). If a combination is in \(S_j\), projection onto the
first augmented flag in the rational direct sum puts its coefficients,
after the unitriangular change of basis, in \(gR\) locally.
The distinguished negative vector is primitive in the middle
quotient \((\Lambda\cap X^\perp)/(\Lambda\cap X)\), so its
coordinate obeys the same condition.
Conversely multiples of \(g\) give zero classes. Hence the subgroup
generated locally has the size of \((R/gR)^p\) globally. Equivalently,
adjoining these marked divisions to \(S_j\) gives an intermediate
lattice of index \(g^{2fp}=g^{mf}\). In particular \(\beta_j\) is
the residual absolute saturation index and
\[
       |gB_j|=|B_j|/|B_j[g]|\le\beta_j .
\]

For \(V=X\) or \(X_j\), write \(L(V)=\Lambda\cap V^\perp\) and
\(V_R=\Lambda\cap V\). Define the enlarged orthogonal lattice
\(A_j\subset W_j^\perp\) by the existence of representatives
\[
 \ell\in L(X),\qquad \ell_j\in L(X_j),\qquad
 x-\ell\in X,\quad x-\ell_j\in X_j,\quad
 \ell-\ell_j\in g\Lambda+X_R+(X_j)_R .
\]
These conditions are \(R\)-linear. The map
\(x\mapsto[\ell-\ell_j]\in B_j\) is independent of representatives,
since two choices on each side differ by an integral flag vector.
Its image belongs to \(gB_j\): subtract the two integral flag
vectors in the congruence and divide by \(g\). The resulting vector
lies in \(\Lambda\cap W_j\).

To identify the kernel, write \(\ell-\ell_j=v+v_j\) in \(S_j\),
with \(v,v_j\) integral in their respective augmented flag spans.
Then
\[
       x-\ell+v=x-\ell_j-v_j
        \ \in\ (X+Ez)\cap(X_j+Ez_j)=0 .
\]
It follows that \(x=\ell-v\in U_j\). Conversely an integral
\(x\in W_j^\perp\) permits the representatives \(\ell=\ell_j=x\).
We have proved
\begin{equation}\label{eq:nonprimitive-one-index}
 I_j=[A_j:U_j]\le\beta_j,\qquad
 d(A_j)I_j^2=d(U_j)I_j\le G_j .
\end{equation}

\proofstep{The integral lattice shared by a matching pair.}
Write \(\eta\) for its integral comparison isometry, carrying the
second flag of the first term to that of the second. The joint
rank is the dimension of the correspondence
\[
 x\in W_1^\perp,\quad y\in W_2^\perp,\qquad
 x-y\in X,\quad \eta x-y\in X_2 .
\]
Each projection is injective since \(W_j^\perp\cap X=0\).
Both target spaces have dimension \(k-1\), so the correspondence
gives a unique \(y\) for each \(x\), and conversely. Projection
modulo \(X\) identifies these rational spaces isometrically.

The marking at the common finite index requires, for every
\(\ell\in L(X)\) and \(\ell_1\in L(X_1)\),
\[
 \ell-\ell_1\in g\Lambda+X_R+(X_1)_R
 \quad\Longleftrightarrow\quad
 \ell-\eta\ell_1\in g\Lambda+X_R+(X_2)_R .
\]
For \(x\in A_1\), keep its first representative \(\ell\) and
use \(\eta\ell_1\) as the second representative of \(y\).
Indeed \(y-\ell\in X\) and
\(y-\eta\ell_1=(y-\eta x)+\eta(x-\ell_1)\in X_2\).
The displayed congruence gives \(y\in A_2\).
The reverse implication and \(\eta^{-1}\) give the converse.
Thus \(A_1,A_2\) have the same projection \(A\) in the
positive definite subspace \(z^\perp\) of \(X^\perp/X\).

In that quotient let \(P_0\) be the intersection of the
projections of \(U_1,U_2\). These are integral sublattices of
the definite middle lattice of the row. The residue map into
\(A/U_1\oplus A/U_2\), with the lattices understood by projection,
now gives
\[
 \begin{aligned}
  {}[A:P_0]&\le I_1I_2,\\
  d(P_0)&\le d(A)I_1I_2
      =\sqrt{d(A_1)I_1^2\,d(A_2)I_2^2}
      \le \sqrt{G_1G_2}.
 \end{aligned}
\]
This uses the level equivalence in both directions. It does not
require either individual frame span to be primitive.

\proofstep{Counting the common sublattices in the row measure.}
Fix reduced labels \(Z_1,Z_2\), write \(H_j=H_{Z_j}\),
\(H=H_1H_2\), \(D=\sqrt{G_1G_2}\), and set \(\lambda=\log_2(2f)\).
The selected prime satisfies \(N=T_u^f\le b^{C_T mf}\le H\),
and Lemma~\ref{lem:reduced-flag-heights} gives \(D\le H\). We may take \(D\ge1\)
whenever there is a pair. For a fixed \(P_0\) with \(d(P_0)\le D\),
Proposition~\ref{prop:integral-flag-frames} applies separately to
the two frames with \(s=1\) and \(Q_j=H_j^2\).
Their classes determine the marked pair class by the uniqueness
argument of Proposition~\ref{prop:flag-pair-counts} (the difference
is a finite-order isometry congruent to the identity modulo \(g\)).
The two bounds in \eqref{eq:flag-over-P} therefore cost together
at most
\[
             9^{2f}H^{8r^2+18r+12}.
\]

If \(P_0\) saturates to determinant \(d_0\), then
\([\overline P_0:P_0]\le D/d_0\). The index count in
Proposition~\ref{prop:primitive-flag-means} bounds all possibilities
over this saturation by \((D/d_0)^{k+\lambda+1}\).
The corank-one part of the same proposition gives the mean count
of saturations under the reciprocal row stabilizer measure.
The two large powers combine as
\[
 (d_0/N)^{k-1}(D/d_0)^{k+\lambda+1}
    =(D/N)^k(N/D)(D/d_0)^{\lambda+2}.
\]

There are at most \(1+D^2\) determinant values since \(d_0^2\)
is an integer. The remaining factors in this average are bounded by
\[
  h_E(4k\sqrt{\mathrm{rd}(E)})^f
        H^{11+3\log_2 f}(1+D^2)(N/D)D^{\lambda+2}.
\]
On the marked range, \((kb)^8\le g\le b^{C_g}\),
\(\mathrm{rd}(E)\le b\), and
\(h_E\le\mathrm{rd}(E)^{f(1+\lambda)}\).
Since \(H_j\ge b^{c^2m^3 f}\) and \(2+\log_2 f\le m\), these
factors and the two frame counts cost at most \(H^{e_0}\) for
\(e_0=C_0m^2\) with \(C_0\) fixed independently of \(C_\Gamma\).

Restoring the remaining corank-one mean factor yields
\[
 \frac{1}{M_u}\sum_{\substack{(t,t')\in\mathcal S\\
                     (t,t')\mapsto(Z_1,Z_2)}}w_{t,t'}
 \le H^{e_0}\left(\frac{e\pi}{k\upsilon}\right)^{kf}
                  (D/N)^k
     = H^{e_0}\sqrt{n_{Z_1}n_{Z_2}},
\]
where \(M_u=m_k(T_u^k-1)^f\)
is the total mass from Proposition~\ref{prop:flag-row-mass}.
The same weighted mean includes every finite coarse row class
of the specified pair subset.

\proofstep{Restoring the archimedean coefficient and the finite norm.}
Use the squared-norm bound with \(D_L=C+6CC_h\) from the proof of
Lemma~\ref{lem:equal-word-generic-corank-one}. Its proof applies to
every positive single-term Gram invariant. For \(d=\dim\pi_u\) it reads
\[
 |a_t|/\sqrt d
 \le H_j^h\,b^{D_Lm^4 f/2}\sqrt{\tau/n_{Z_j}}\,
                         \sigma_{Z_j}^{C_2m^2/2-a}.
\]

Thus the square roots of the two \(n_{Z_j}\) cancel on this very
fibre. After \(H_j\le P\sigma_{Z_j}^4\), the exponent of each
\(\sigma_{Z_j}\) is \(4(e_0+h)+C_2m^2/2-a\).
The actual order \(a=\lceil pa_0\rceil/p\) of
Lemma~\ref{lem:row-moment-budget} satisfies \(a\ge4ck\).
Since \(k\ge m^5(\log b)x^2\), eventually
\[
        a\ge4(e_0+h)+|C_2|m^2/2+12p^2 .
\]

Lemma~\ref{lem:gram-label-sum} now bounds the two independent
label sums by \(\Phi^2\). With
\(\Gamma_1M_ud=d_u\) and \(\tau=(1+\delta)d_u/D_*\),
the resulting bound is
\[
     (1+\delta)P^{2(e_0+h)}\Phi^2 b^{D_Lm^4f}\,\frac{d_u^2}{D_*}.
\]
Indeed the weights just counted are the finite row weights, while
the two coefficient norms use the real orbital Haar of
Proposition~\ref{prop:two-row-amplitude-bounds}.
Their only remaining finite multiplier is \(\Gamma_1\).
Here \(\log P=O(m^3f\log b)\), \(e_0+h=O(m^2)\), and
\(\Phi=(C_*^fP^3)^{p^2}\), proving \eqref{eq:nonprimitive-diagonal}.
\end{proof}
